\chapter{การออกแบบและปรับแต่งประสิทธิภาพ SQL บน Google BigQuery}

Google BigQuery เป็นระบบคลังข้อมูลระดับองค์กรที่ใช้สถาปัตยกรรมแยกส่วนการคำนวณออกจากพื้นที่จัดเก็บข้อมูล (Decoupled Compute and Storage) ทำให้สามารถขยายขีดความสามารถในการสืบค้นข้อมูลขนาดมหาศาลได้อย่างรวดเร็ว

\section{สถาปัตยกรรมกระจายการประมวลผล Dremel และ Colossus}

ระบบ Dremel จัดแบ่งงานคิวรีออกเป็นโครงสร้างต้นไม้ (Execution Tree) โดยมี Slot เป็นหน่วยประมวลผลเสมือน การลดปริมาณข้อมูลที่อ่านจากดิสก์ Colossus จึงเป็นปัจจัยชี้ขาดต่อทั้งความเร็วและค่าใช้จ่าย

\section{เทคนิค Partitioning และ Clustering}

การแบ่งพาร์ติชัน (Partitioning) ตามคอลัมน์วันที่หรือเวลา ช่วยให้คำสั่งคิวรีทำการ Pruning ข้อมูลเฉพาะช่วงที่ต้องการ:

\begin{equation}
\text{Data Scanned Ratio} = \frac{\sum \text{Bytes of Filtered Partitions}}{\sum \text{Bytes of Entire Table}} \ll 1
\end{equation}

เมื่อใช้ร่วมกับการจัดกลุ่มข้อมูล (Clustering) ตามคอลัมน์ที่มีการกรองบ่อยครั้ง เช่น รหัสสถานีตรวจวัด หรือรหัสประเภทเซนเซอร์ BigQuery จะสามารถข้ามการอ่านบล็อกข้อมูลที่ไม่เกี่ยวข้องได้มากกว่าร้อยละ 80

\section{กลยุทธ์การเขียนคำสั่ง SQL เพื่อหลีกเลี่ยง Data Shuffling}

การเชื่อมโยงตารางขนาดใหญ่ (Large Joins) ควรกำหนดให้ตารางที่มีขนาดเล็กกว่าอยู่ทางฝั่งขวา เพื่อให้ระบบทำการบรอดแคสต์ (Broadcast Hash Join) แทนที่จะกระจายข้อมูลข้ามโหนด ซึ่งช่วยลดเวลาการประมวลผลและป้องกันข้อผิดพลาดจากทรัพยากร Slot ไม่เพียงพอ
